\section{Định nghĩa tổ hợp}

Công thức cơ bản:

\[
C(n,k) = \binom{n}{k} = \frac{n!}{k!(n-k)!}
\]

Ý nghĩa: số cách chọn $k$ phần tử từ $n$ phần tử.

\section{Chọn phương pháp phù hợp}

\begin{center}
\begin{tabular}{|c|c|}
\hline
Trường hợp & Phương pháp \\
\hline
$n \le 20$ & Công thức trực tiếp \\
\hline
$n \le 2000$ & Tam giác Pascal \\
\hline
$n \le 10^6$, mod prime & Giai thừa + nghịch đảo \\
\hline
$n$ rất lớn, mod nhỏ & Định lý Lucas \\
\hline
$n$ lớn, $k$ nhỏ, mod prime & Công thức trực tiếp + nghịch đảo \\
\hline
\end{tabular}
\end{center}

\newpage

\section{Tính trực tiếp (Multiplicative Formula)}

Đây là cách tính tổ hợp dựa trực tiếp vào công thức:

\[
C(n,k) = \frac{n!}{k!(n-k)!}
\]

Tuy nhiên nếu tính trực tiếp bằng giai thừa thì dễ gây tràn số và tốn nhiều phép tính.  
Ta biến đổi công thức thành dạng nhân tuần tự:

\[
C(n,k) =
\frac{n(n-1)(n-2)\dots(n-k+1)}{k!}
\]

hay viết lại:

\[
C(n,k) =
\prod_{i=1}^{k} \frac{n-i+1}{i}
\]

Công thức này cho phép ta tính dần từng bước, mỗi bước nhân thêm một số ở tử và chia cho một số ở mẫu.

\subsection*{Ý tưởng}

Ban đầu đặt:

\[
res = 1
\]

Sau đó lặp từ $i = 1$ đến $k$:

\[
res = res \times \frac{n-i+1}{i}
\]

Nhờ tính chất của tổ hợp, tại mỗi bước giá trị thu được luôn là số nguyên nên phép chia luôn hợp lệ.

\subsection*{Độ phức tạp}

\begin{itemize}
\item Thời gian: $O(k)$
\item Bộ nhớ: $O(1)$
\end{itemize}

\subsection*{Khi nào nên dùng}

Phương pháp này phù hợp khi:

\begin{itemize}
\item $n$ không quá lớn
\item $k$ nhỏ
\item Không yêu cầu tính theo modulo
\item Chỉ cần tính một vài giá trị tổ hợp
\end{itemize}

Nếu $n$ lớn hoặc cần tính nhiều truy vấn, nên sử dụng phương pháp \textit{precompute factorial}.

\subsection*{Cài đặt}

\begin{tcolorbox}
\begin{minted}{cpp}
long long C(int n, int k) {
    if (k < 0 || k > n) return 0; // Điều kiện ràng buộc
    long long res = 1;
    for(int i = 1; i <= k; i++)
        res = res * (n - i + 1) / i;
    return res;
}
\end{minted}
\end{tcolorbox}

\subsection*{Tối ưu nhỏ}

Do tính chất đối xứng của tổ hợp:

\[
C(n,k) = C(n,n-k)
\]

ta có thể giảm số vòng lặp bằng cách:

\[
k = \min(k, n-k)
\]

Điều này giúp thuật toán chạy nhanh hơn khi $k$ gần $n$.

\subsection*{Ví dụ}

Tính:

\[
C(5,2)
\]

Ta có:

\[
C(5,2) =
\frac{5 \times 4}{2 \times 1}
= 10
\]

\newpage

\section{Pascal Triangle (Dynamic Programming)}

Phương pháp này dựa trên công thức truy hồi nổi tiếng của tổ hợp:

\[
C(n,k) = C(n-1,k-1) + C(n-1,k)
\]

Công thức này xuất phát từ việc phân chia các trường hợp khi chọn $k$ phần tử từ $n$ phần tử:

\begin{itemize}
\item Chọn phần tử thứ $n$: khi đó cần chọn thêm $k-1$ phần tử từ $n-1$ phần tử còn lại $\Rightarrow C(n-1,k-1)$.
\item Không chọn phần tử thứ $n$: khi đó chọn $k$ phần tử từ $n-1$ phần tử $\Rightarrow C(n-1,k)$.
\end{itemize}

Do đó:

\[
C(n,k) = C(n-1,k-1) + C(n-1,k)
\]

\subsection*{Điều kiện biên}

Ta có các giá trị đặc biệt:

\[
C(n,0) = 1
\]

\[
C(n,n) = 1
\]

vì:

\begin{itemize}
\item Chọn 0 phần tử luôn có đúng 1 cách
\item Chọn toàn bộ $n$ phần tử cũng chỉ có 1 cách
\end{itemize}

\subsection*{Tam giác Pascal}

Các giá trị $C(n,k)$ tạo thành một cấu trúc gọi là \textbf{Tam giác Pascal}:

\[
\begin{matrix}
&&1\\
&1&&1\\
1&&2&&1\\
1&&3&&3&&1\\
1&&4&&6&&4&&1
\end{matrix}
\]

Mỗi số bên trong được tính bằng tổng của hai số phía trên.

\subsection*{Ý tưởng thuật toán}

Ta xây dựng bảng $C[n][k]$ bằng quy hoạch động:

\begin{enumerate}
\item Khởi tạo các giá trị biên $C(i,0)=1$ và $C(i,i)=1$
\item Với mỗi $i$ từ $1$ đến $n$
\item Tính các giá trị bên trong bằng công thức truy hồi
\end{enumerate}

\subsection*{Cài đặt}

\begin{tcolorbox}
\begin{minted}{cpp}
const int MAXN = 2000;
long long C[MAXN][MAXN];

void build() {
    for(int i = 0; i < MAXN; i++) {
        C[i][0] = C[i][i] = 1;
        for(int j = 1; j < i; j++)
            C[i][j] = C[i-1][j-1] + C[i-1][j];
    }
}
\end{minted}
\end{tcolorbox}

Sau khi gọi hàm \texttt{build()}, giá trị tổ hợp có thể truy cập trực tiếp:

\[
C(n,k) = C[n][k]
\]

\subsection*{Độ phức tạp}

\begin{itemize}
\item Thời gian tiền xử lý: $O(n^2)$
\item Truy vấn mỗi tổ hợp: $O(1)$
\item Bộ nhớ: $O(n^2)$
\end{itemize}

\subsection*{Tính tổ hợp theo modulo}

Trong lập trình thi đấu, giá trị tổ hợp thường rất lớn nên cần tính theo modulo:

\[
C(n,k) = (C(n-1,k-1) + C(n-1,k)) \bmod m
\]

Ta chỉ cần thêm phép modulo vào công thức:

\begin{tcolorbox}
\begin{minted}{cpp}
const int MAXN = 2000;
const long long MOD = 1000000007;

long long C[MAXN][MAXN];

void build() {
    for(int i = 0; i < MAXN; i++) {
        C[i][0] = C[i][i] = 1;
        for(int j = 1; j < i; j++)
            C[i][j] = (C[i-1][j-1] + C[i-1][j]) % MOD;
    }
}
\end{minted}
\end{tcolorbox}

\subsection*{Khi nào nên dùng}

Phương pháp Pascal Triangle phù hợp khi:

\begin{itemize}
\item $n$ không quá lớn (thường $n \le 2000$)
\item Cần tính nhiều giá trị $C(n,k)$
\item Modulo không phải số nguyên tố
\end{itemize}

\subsection*{Nhược điểm}

\begin{itemize}
\item Tốn bộ nhớ $O(n^2)$
\item Không phù hợp khi $n$ rất lớn ($n > 10^5$)
\end{itemize}

Trong các bài toán lập trình thi đấu với $n$ lớn, phương pháp \textit{factorial + modular inverse} thường được sử dụng thay thế.

\newpage

\section{Precompute factorial + inverse (Phổ biến nhất)}

Đây là phương pháp được sử dụng nhiều nhất trong lập trình thi đấu khi:

\begin{itemize}
\item $n$ có thể lớn (tới $10^6$ hoặc hơn)
\item Có nhiều truy vấn tính $C(n,k)$
\item Kết quả cần tính theo modulo
\item Modulo là số nguyên tố (thường là $10^9+7$ hoặc $998244353$)
\end{itemize}

\subsection*{Công thức}

Tổ hợp được định nghĩa bởi:

\[
C(n,k) = \frac{n!}{k!(n-k)!}
\]

Khi tính theo modulo $p$:

\[
C(n,k) = n! \times (k!)^{-1} \times ((n-k)!)^{-1} \pmod p
\]

Vấn đề chính là tính \textbf{nghịch đảo modular} của giai thừa.

\subsection*{Định lý Fermat nhỏ}

Nếu $p$ là số nguyên tố và $a$ không chia hết cho $p$ thì:

\[
a^{-1} \equiv a^{p-2} \pmod p
\]

Do đó:

\[
(k!)^{-1} = (k!)^{p-2} \pmod p
\]

Ta có thể dùng \textbf{binary exponentiation} để tính lũy thừa nhanh.

\subsection*{Ý tưởng thuật toán}

Thay vì tính tổ hợp mỗi lần từ đầu, ta tiền xử lý hai mảng:

\begin{itemize}
\item \textbf{fact[i]} = $i! \bmod p$
\item \textbf{invfact[i]} = $(i!)^{-1} \bmod p$
\end{itemize}

Sau khi tiền xử lý xong, ta có:

\[
C(n,k) =
fact[n] \times invfact[k] \times invfact[n-k] \pmod p
\]

Nhờ vậy mỗi truy vấn tổ hợp chỉ mất $O(1)$.

\subsection*{Các cách tính nghịch đảo giai thừa}

Có nhiều cách để tính mảng nghịch đảo giai thừa. Dưới đây là hai phương pháp phổ biến.

\subsubsection*{Cách 1: Tính từ $(N!)^{-1}$ rồi suy ra các giá trị nhỏ hơn}

Ta sử dụng định lý Fermat để tính:

\[
(N!)^{-1} = (N!)^{p-2} \pmod p
\]

Sau đó suy ra các giá trị còn lại bằng công thức:

\[
(i-1)!^{-1} = i!^{-1} \times i \pmod p
\]

Ý tưởng:

\begin{itemize}
\item Tính toàn bộ mảng $fact$
\item Tính $invfact[N]$
\item Duyệt ngược để tính các giá trị còn lại
\end{itemize}

\begin{tcolorbox}
\begin{minted}{cpp}
fact[0] = 1;
for(int i = 1; i <= N; i++)
    fact[i] = fact[i-1] * i % MOD;

invfact[N] = power(fact[N], MOD-2);

for(int i = N; i >= 1; i--)
    invfact[i-1] = invfact[i] * i % MOD;
\end{minted}
\end{tcolorbox}

\textbf{Ưu điểm}

\begin{itemize}
\item Chỉ cần một phép lũy thừa modulo
\item Thời gian tiền xử lý $O(N)$
\end{itemize}

Đây là cách được sử dụng phổ biến nhất trong lập trình thi đấu.

\subsubsection*{Cách 2: Tính nghịch đảo của từng số rồi nhân dần}

Thay vì tính nghịch đảo của giai thừa trực tiếp, ta tính trước nghịch đảo của các số:

\[
inv[i] = i^{-1} \pmod p
\]

Có công thức quy hoạch động:

\[
inv[i] = p - \left(\frac{p}{i}\right) \times inv[p \bmod i] \pmod p
\]

Sau đó:

\[
invfact[i] = invfact[i-1] \times inv[i] \pmod p
\]

\begin{tcolorbox}
\begin{minted}{cpp}
inv[1] = 1;
for(int i = 2; i <= N; i++)
    inv[i] = MOD - MOD/i * inv[MOD % i] % MOD;

fact[0] = invfact[0] = 1;

for(int i = 1; i <= N; i++){
    fact[i] = fact[i-1] * i % MOD;
    invfact[i] = invfact[i-1] * inv[i] % MOD;
}
\end{minted}
\end{tcolorbox}

\textbf{Ưu điểm}

\begin{itemize}
\item Không cần phép lũy thừa modulo
\item Tính được toàn bộ nghịch đảo $1..N$ trong $O(N)$
\end{itemize}

\textbf{Nhược điểm}

\begin{itemize}
\item Phải lưu thêm mảng $inv$
\end{itemize}

\subsection*{So sánh}

\begin{center}
\begin{tabular}{|c|c|c|}
\hline
Phương pháp & Lũy thừa modulo & Bộ nhớ \\
\hline
Từ $(N!)^{-1}$ & 1 lần & $fact$, $invfact$ \\
\hline
Từ $1/i$ & Không cần & $fact$, $invfact$, $inv$ \\
\hline
\end{tabular}
\end{center}

\subsection*{Bước tiền xử lý}

\begin{enumerate}
\item Tính mảng $fact$
\item Tính $invfact[n]$ bằng Fermat
\item Tính ngược các $invfact$
\end{enumerate}

\subsection*{Cài đặt}

\begin{tcolorbox}
\begin{minted}{cpp}
long long C(int n, int k){
    if(k < 0 || k > n) return 0;
    return fact[n] * invfact[k] % MOD * invfact[n - k] % MOD;
}
\end{minted}
\end{tcolorbox}

\subsection*{Độ phức tạp}

\begin{itemize}
\item Tiền xử lý giai thừa: $O(n)$
\item Tính nghịch đảo: $O(\log p)$
\item Truy vấn mỗi tổ hợp: $O(1)$
\item Bộ nhớ: $O(n)$
\end{itemize}

\subsection*{Ưu điểm}

\begin{itemize}
\item Tính tổ hợp cực nhanh sau khi tiền xử lý
\item Phù hợp cho bài toán có nhiều truy vấn
\item Có thể xử lý $n$ rất lớn (tới $10^6$ hoặc hơn)
\end{itemize}

\subsection*{Nhược điểm}

\begin{itemize}
\item Chỉ áp dụng khi modulo là số nguyên tố
\item Cần bộ nhớ để lưu các mảng giai thừa
\end{itemize}

\subsection*{Ví dụ}

Giả sử cần tính:

\[
C(10,3)
\]

Ta có:

\[
C(10,3)=
\frac{10!}{3!\,7!}
=120
\]

Trong chương trình:

\[
C(10,3) =
fact[10] \times invfact[3] \times invfact[7]
\]

\subsection{Trường hợp $n$ rất lớn nhưng $k$ nhỏ}

Trong nhiều bài toán lập trình thi đấu, ta gặp trường hợp:

\begin{itemize}
\item $n$ rất lớn (có thể tới $10^{12}$ hoặc $10^{18}$)
\item $k$ nhỏ (thường $\le 10^6$ hoặc nhỏ hơn)
\end{itemize}

Trong trường hợp này, việc tiền xử lý $n!$ là không khả thi. Ta có thể sử dụng trực tiếp công thức nhân:

\[
C(n,k) =
\frac{n(n-1)(n-2)\dots(n-k+1)}{k!}
\]

hay viết dưới dạng tích:

\[
C(n,k) =
\prod_{i=1}^{k}
\frac{n-i+1}{i}
\]

\subsubsection*{Ý tưởng}

Ta nhân dần các số ở tử số và chia cho mẫu số tương ứng.

Để tránh tràn số và đảm bảo tính đúng khi làm việc với modulo, ta thường:

\begin{itemize}
\item Tiền xử lý $k!$
\item Tính nghịch đảo của $k!$
\item Nhân tử số $k$ số liên tiếp
\end{itemize}

\subsubsection*{Cài đặt}

\begin{tcolorbox}
\begin{minted}{cpp}
long long C(long long n, int k){
    if(k < 0 || k > n) return 0;
    long long res = 1;

    for(int i = 1; i <= k; i++){
        res = res * (n - i + 1) % MOD;
        res = res * inv[i] % MOD;
    }

    return res;
}
\end{minted}
\end{tcolorbox}

Trong đó:

\begin{itemize}
\item $inv[i] = i^{-1} \pmod{MOD}$
\item mảng $inv$ có thể được tiền xử lý trong $O(k)$
\end{itemize}

\subsubsection*{Ví dụ}

Tính:

\[
C(10^{12},3)
\]

Ta có:

\[
C(n,3)=\frac{n(n-1)(n-2)}{6}
\]

Chỉ cần nhân ba số lớn thay vì tính $n!$.

\subsubsection*{Độ phức tạp}

\begin{itemize}
\item Thời gian: $O(k)$
\item Bộ nhớ: $O(k)$
\end{itemize}

\subsubsection*{Khi nào nên dùng}

Phương pháp này rất hiệu quả khi:

\begin{itemize}
\item $n$ cực lớn
\item $k$ nhỏ
\item modulo là số nguyên tố
\end{itemize}

Trong nhiều bài toán Codeforces hoặc ICPC, đây là phương pháp đơn giản và hiệu quả nhất cho các tổ hợp có tham số $n$ rất lớn.

\newpage

\section{Lucas Theorem}

Trong một số bài toán lập trình thi đấu, giá trị $n$ có thể rất lớn (thậm chí tới $10^{18}$ hoặc hơn), khiến việc tiền xử lý giai thừa trở nên không khả thi. Khi đó ta có thể sử dụng \textbf{Lucas Theorem} để tính tổ hợp theo modulo một số nguyên tố nhỏ.

\subsection*{Phát biểu định lý}

Giả sử $p$ là số nguyên tố. Khi đó:

\[
\binom{n}{k}
\equiv
\prod_{i=0}^{m}
\binom{n_i}{k_i}
\pmod p
\]

trong đó:

\begin{itemize}
\item $n_i, k_i$ là các chữ số của $n$ và $k$ khi viết trong hệ cơ số $p$
\item $n = n_0 + n_1 p + n_2 p^2 + \dots$
\item $k = k_0 + k_1 p + k_2 p^2 + \dots$
\end{itemize}

Nếu tồn tại $k_i > n_i$ thì:

\[
\binom{n}{k} \equiv 0 \pmod p
\]

\subsection*{Ý tưởng}

Thay vì tính trực tiếp $C(n,k)$ với $n$ rất lớn, ta:

\begin{enumerate}
\item Phân tích $n$ và $k$ theo hệ cơ số $p$
\item Tính từng tổ hợp nhỏ $C(n_i,k_i)$
\item Nhân tất cả lại theo modulo $p$
\end{enumerate}

Các giá trị $C(n_i,k_i)$ đều có:

\[
0 \le n_i, k_i < p
\]

nên có thể tính bằng các phương pháp thông thường (Pascal hoặc factorial).

\subsection*{Triển khai đệ quy}

Lucas theorem thường được cài đặt bằng đệ quy:

\[
C(n,k) =
C(n \bmod p, k \bmod p)
\times
C(\lfloor n/p \rfloor, \lfloor k/p \rfloor)
\]

\subsection*{Cài đặt}

Giả sử ta đã có hàm $C(n,k)$ tính tổ hợp với $n < p$.

\begin{tcolorbox}
\begin{minted}{cpp}
long long lucas(long long n, long long k){
    if(k == 0) return 1;

    return C(n % MOD, k % MOD) *
           lucas(n / MOD, k / MOD) % MOD;
}
\end{minted}
\end{tcolorbox}

\subsection*{Ví dụ}

Tính:

\[
C(1000,300) \pmod 13
\]

Ta viết:

\[
1000 = 5 \cdot 13^2 + 11 \cdot 13 + 12
\]

\[
300 = 1 \cdot 13^2 + 10 \cdot 13 + 1
\]

Theo Lucas theorem:

\[
C(1000,300)
\equiv
C(12,1) \times C(11,10) \times C(5,1)
\pmod{13}
\]

Các giá trị này đều nhỏ hơn $13$ nên dễ dàng tính được.

\subsection*{Độ phức tạp}

Giả sử $n$ có $d$ chữ số trong hệ cơ số $p$:

\begin{itemize}
\item Số bước đệ quy: $O(\log_p n)$
\item Mỗi bước tính một tổ hợp nhỏ
\end{itemize}

Do đó tổng độ phức tạp:

\[
O(\log_p n)
\]

\subsection*{Khi nào nên dùng}

Lucas theorem rất hữu ích khi:

\begin{itemize}
\item $n$ rất lớn (có thể tới $10^{18}$)
\item modulo là số nguyên tố nhỏ
\item không thể tiền xử lý giai thừa tới $n$
\end{itemize}

\subsection*{Nhược điểm}

\begin{itemize}
\item Chỉ áp dụng khi modulo là số nguyên tố
\item Không hiệu quả nếu modulo quá lớn
\end{itemize}

Trong thực tế, Lucas theorem thường được dùng khi:

\begin{itemize}
\item $p \le 10^6$
\item $n$ rất lớn
\item số truy vấn ít
\end{itemize}

\newpage

\section{Kết luận}

Trong hầu hết các bài toán lập trình thi đấu:

\begin{itemize}
\item Precompute factorial + inverse là phương pháp phổ biến nhất.
\item Lucas theorem dùng khi $n$ rất lớn.
\end{itemize}
